\section{Conclusions and Future Directions}
\label{sec:discussion}
\subsection{Summary of the Results}
\label{sec:farm-summary}
% Cells computed from analysis/rev4_farm_tables.json (tables.*.ranking_order, methods.*.rank, psovns_vs.outcome; Type I: tab:farm-I-stats avg_rank, vs_psovns p_wilcoxon_holm).
\begin{table}[!htb]
\centering
\caption{Summary of the nine methods by wind farm type: rank-1 method, ranks of PSO-VNS and PSO, and significant differences of PSO-VNS against the other methods.}
\label{tab:farm-summary}
\footnotesize\setlength{\tabcolsep}{3pt}
\begin{tabular}{llcccc}
\toprule
Farm type & Budget & Rank-1 method & PSO-VNS & PSO & W/L \\
\midrule
I: benchmark (68 cases) & 6{,}030 & PSO-VNS & 1 & 2 & 6/0 \\
II: Horns Rev~1 & 6{,}030 & PSO-VNS & 1 & 3 & 7/0 \\
 & 30{,}030 & GA & 2 & 3 & 7/1 \\
III: Lillgrund & 6{,}030 & PSO-VNS & 1 & 7.5 & 6/1 \\
 & 30{,}030 & PSO-VNS & 1 & 6 & 7/0 \\
IV: IEA37, 16 turbines & 6{,}030 & PSO-VNS & 1 & 3 & 8/0 \\
 & 30{,}030 & PSO-VNS & 1 & 3 & 7/0 \\
IV: IEA37, 36 turbines & 6{,}030 & GA & 3 & 4 & 5/1 \\
 & 30{,}030 & GA & 2 & 3 & 6/1 \\
\bottomrule
\end{tabular}
\par\vspace{2pt}\parbox{\columnwidth}{\scriptsize Rank: Type~I, position by average rank (Table~\ref{tab:farm-I-stats}); Types~II--IV, feasibility-aware rank of the mean AEP (Tables~\ref{tab:farm-II}--\ref{tab:farm-IV}; 7.5: tied, no feasible run). W/L: methods to which PSO-VNS is significantly better / worse among its 8 comparisons (Holm-adjusted, $\alpha=0.05$; Type~I: case-mean Wilcoxon test, Types~II--IV: seed-paired Wilcoxon test).}
\end{table}

PSO-VNS ranks first in six of the nine settings of Table~\ref{tab:farm-summary} and second or third in the other three, in which GA ranks first: the Horns Rev~1 block at 30,030 evaluations and both IEA37 budgets with 36 turbines. PSO-VNS ranks ahead of PSO in every setting and is significantly better than PSO in every site setting, but not on the benchmark, where the two are practically equivalent at the seed level (Section~\ref{sec:psovns-pso}). PSO-VNS is significantly worse than one method in four settings (GA three times; MS-SLSQP on the Lillgrund block at 6,030 evaluations, where it is feasible in all runs) and never worse than more than one.

\subsection{Where PSO-VNS Helps and Where It Does Not}
\label{sec:disc-where}
PSO-VNS
\begin{itemize}
\item reaches feasibility more reliably than PSO on the two site blocks,
\item ranks best on the six largest cases at all three budgets (exploratory) and
\item has a lower mean wake loss than PSO in all ten Data Set~II cases with $N\ge10$ (post hoc), also with 1$^\circ$ bins.
\end{itemize}
It does not help
\begin{itemize}
\item from feasible starts,
\item in the densest Data Set~I layouts,
\item on IEA37 with 36 turbines (GA ranks first) or against exact-gradient SLSQP,
\item at the larger Horns Rev~1 budget (GA leads),
\item in reliability on Lillgrund at 6,030 evaluations (feasible runs: MS-SLSQP 30, PSO-VNS 17 of 30) or
\item on the 15$^\circ$ benchmark average by more than a few hundredths of a percentage point,
\end{itemize}
and its Lillgrund AEP lead over the runners-up is not significant with 1$^\circ$ bins.
% CHECK-EXTRA [X52]: Horns Rev 30/30 vs 19/30, McNemar p 9.8e-4 (one site, random starts); CHECK-FINAL [C08]/[C20]; generated \NFeasPSO = 100 (R5 O8); [C29]: six largest cases best at all budgets (exploratory)
% CHECK-FINAL [C32] + CHECK-EXTRA [X19]: DS II N >= 10, all-family Holm; CHECK-DIR [F14]: survives 1 deg; CHECK-EXTRA [X43]: all \NXHetDsINPSO DS I cases beyond m favour PSO; [X25]: mean gain < 0.05 % AEP; CHECK-FINAL [C28]/[C40]
% CHECK (manual, 2026-10-04): Lillgrund 17/30 vs 0/30, 30/30 vs 19/30; at 1 deg PSO-VNS highest mean at 30,030 but p_Holm 0.16 / 0.58 / 0.58 vs SSA-VNS / VNS / MS-SLSQP, the runners-up (tab:S-sa-lg-eval); GA first at 30,030 on Horns Rev 1 with 5-deg, 1-deg (+0.48, p 2.8e-4) and PyWake NOJ (+0.43, p 4.2e-4) (tab:S-sa-hr-eval; analysis/rev3_sites.json)

\subsection{Implications for Wind Farm Layout Practice}
\label{sec:disc-practice}
The direction discretization is the main caveat for practice: with 1$^\circ$ bins the benchmark wake loss rises by 0.40--1.04~pp, no optimized Horns Rev~1 layout exceeds the installed block (Section~\ref{sec:hornsrev-site}), and optimizing with 1$^\circ$ instead of 15$^\circ$ bins lowers the 1$^\circ$ wake loss by up to 1.1~pp (Section~\ref{sec:direction}). Gains of layouts optimized under a coarse rose should be verified, and preferably obtained, at a fine resolution before they are credited to an optimizer.
% CHECK-DIR [F03]: every method's loss rises by \NFLossRiseMin to \NFLossRiseMax pp; [F21]: 0 of 961 feasible runs (ten methods and GA) above the installed block at 1 deg (22 at 5 deg; ten methods only: 0 of 901, 16); CHECK-FINAL [C62]: no \NHR...Above macro used; CHECK (manual, 2026-10-04): direct 1-deg optimization lowers the 1-deg loss by 0.33-1.12 pp on the same seeds (analysis/rev3_fine.json)

The Horns Rev~1 block (16 of 80 turbines, $4D$ instead of $7D$ spacing, no cables, loads or navigation) tests optimizers, not the installed design (Section~\ref{sec:hornsrev-site}). On IEA37, exact gradients gave a higher AEP than every metaheuristic at the same charged budget (Section~\ref{sec:iea37}); for smooth, differentiable wake models, our results thus favor gradient-based optimization, possibly started from a swarm. The metaheuristic comparisons of this paper are relevant where the objective is discontinuous, as with the top-hat wake of the benchmark, or gradients are unavailable.

\subsection{Recommendations for Benchmarking}
\label{sec:disc-recs}
Table~\ref{tab:recommendations} condenses the protocol into a checklist~\cite{BartzBeielstein2020,LaTorre2021}.
\begin{table}[!t]
\centering
\caption{Checklist for benchmarking metaheuristics on the WFLOP.}
\label{tab:recommendations}
\footnotesize
\begin{tabular}{@{}>{\raggedright\arraybackslash}p{0.02\textwidth}>{\raggedright\arraybackslash}p{0.36\textwidth}>{\raggedright\arraybackslash}p{0.27\textwidth}>{\raggedright\arraybackslash}p{0.25\textwidth}@{}}
\toprule
& Item & Rationale in this study & What we did (where) \\
\midrule
1 & Equal model-query budgets (incl.\ initial population, gradients), elapsed time reported separately; seed-paired runs & Iterations differ in cost (LX-SSA: $2N_p$) & $B=6{,}030$; seeds 1--30 (Section~\ref{sec:setup}) \\
2 & Report feasibility rates, conditional quality and an all-run paired outcome separately & Conditional means can favor less reliable methods & Qualification threshold (Section~\ref{sec:setup-stats}); all-run endpoint on all 68 cases and Lillgrund (Section~\ref{sec:archive-reliability}) \\
3 & Paired, multiplicity-corrected tests; robustness to thresholds and clusters & Rank post hoc tests depend on the pool; cases share farms & Holm, six thresholds, six clusters (Section~\ref{sec:setup}) \\
4 & Equivalence at a stated level, margin and evaluator; name the primary questions; report minimal margins & Verdicts change with the level & TOST at three levels (Section~\ref{sec:psovns-pso}) \\
5 & Ablate hybrids against the local method alone and an equal-cost in-site random control & A square control flatters salp-swarm phases & VNS, RS-VNS, RSD-VNS (Table~\ref{tab:ablation}) \\
6 & Vary budget and initialization before generalizing a ranking & The leader changes from feasible starts & Three budgets, feasible starts (Table~\ref{tab:feasbudget}) \\
7 & Re-evaluate at a finer rose and other models, and re-optimize, before claiming a robust ranking; publish versioned code and data & Layouts exploit coarse bins & 1$^\circ$ re-evaluation and direct 1$^\circ$ optimization (8 cases), Gaussian, PyWake (Sections~\ref{sec:hornsrev-site} and~\ref{sec:robustness}). Code and data: versioned archive (release 1.0.0), available on request \\
8 & Test the conclusions under an alternative constraint handling (normalized violations, feasibility rules, projection) & The order changed with box clipping and Deb's rules & Deb's rules with clipping or projection, six cases (Section~\ref{sec:boundary}) \\
\bottomrule
\end{tabular}
\end{table}
% CHECK-DIAG [D01]/[D02]: instrumented runs reproduce the stored runs; CHECK-EXTRA [X02]-[X09], [X11]/[X14], [X18] (item 3); [X41]: minimal margin grows with the level (item 4); CHECK-FINAL [C57]/[C60] (item 5); [C28] (item 6); CHECK-DIR [F03]/[F07] (item 7); CHECK (manual, 2026-10-04): item 2 all-run endpoint on 68 cases (tab:S-sa-allrun); item 7 direct 1-deg optimization on 8 cases (rev3_fine.json); item 8 constraint study, GA best case-level rank under Deb + clipping (rev3_constraint.json)

\subsection{Limitations}
\label{sec:limits}
Each threat below names the claims it limits; unless a check is cited, the threat is disclosed rather than resolved.

\emph{Internal validity.} PSO-VNS repeats PSO exactly up to the switch, and the ablation variants differ from it only in Phase~1 (controls: 3,015 versus 3,030 evaluations); the code reproduces the archived runs of the original implementation and the stored runs (Section~\ref{sec:validation}). Configuration and constraint handling remain threats:
\begin{itemize}
\item DE is untuned (its low feasibility, 71.4\%, depends on the repair; Section~\ref{sec:boundary});
\item MS-SLSQP uses finite differences on a discontinuous objective (exact gradients: Section~\ref{sec:iea37});
\item the penalty adds boundary violations in m$^2$ to spacing violations in m ($\mu=10^{10}$), box clipping leaves the circle to the penalty, and the conclusions partly depend on this choice (Section~\ref{sec:boundary}); the weaknesses observed for SSA, LX-SSA and DE therefore characterize these implementations under this penalty and repair, not their algorithm families;
\item the SSA--LX-SSA contrast of the main comparison confounds the Laplace step with the doubled cost per iteration; the ablation that separates them (Section~\ref{S-sec:S-xp-lap}) covers only the twelve cases of the split study, on which none of its Holm-adjusted case-mean tests against LX-SSA is significant.
\end{itemize}
% CHECK (manual, see 05_setup): 2 of 24 Mann-Whitney tests reject; CHECK-DIAG [D01]/[D02]: instrumented runs reproduce the stored runs; CHECK-FINAL [C61]: controls switch at 3,015, swarms at 3,030; generated \NFeasDE ("depends on the repair": DE feasibility 33.3 % -> 80.0 % with Deb + projection, rev3_constraint.json); CHECK-FINAL [C17]: LX-SSA as published; CHECK (manual, 2026-10-04): package re-execution, repro_package_build/verification/verification_summary.txt (see 05_setup sec:validation)

\emph{Construct validity.} At 6,030 evaluations the ``VNS'' phase is essentially one truncated compass-search descent, so the ablation concerns this local search, not VNS in general. The objective is a benchmark value, not farm quality; final quality at a fixed budget is the only measure tested; and the qualification-then-conditional-quality ranking limits, but does not remove, survivorship bias.
% CHECK-DIAG [D08]-[D10]: truncated compass search; CHECK-FINAL [C21]: cubic curve lowers the power, ranking intact

\emph{External validity.} Except in the direct 1$^\circ$ study on eight cases, the search used coarse roses (15$^\circ$ benchmark bins, 5$^\circ$ on the sites), which optimized layouts exploit. Further limits:
\begin{itemize}
\item the top-hat Jensen wake of the benchmark implementation, with its non-physical upstream deficit (Section~\ref{sec:model}), and its linear power curve (Section~\ref{sec:powermodel});
\item circular farms, a $4D$ spacing and a prescribed $N$, with random starts and 6,030 evaluations unless stated: rankings of a small budget;
\item measured wind climates enter only through two 16-turbine blocks, Horns Rev~1 and Lillgrund (seven months of measurements, minimum spacing $3D$), on which the feasibility advantage of PSO-VNS rests and whose AEP rankings partly depend on the evaluator (Sections~\ref{sec:hornsrev-site} and~\ref{sec:lillgrund}); the 64-turbine IEA37 scenario was not run;
\item the pool lacks CMA-ES~\cite{HansenOstermeier2001}, L-SHADE~\cite{Tanabe2014}, adaptive GAs, ACO~\cite{Eroglu2012} and BBO~\cite{Bansal2017}; DE and a real-coded GA are its only evolutionary methods, so ``first'' means first among the nine methods compared;
\item the exact-gradient baselines were run only on IEA37 Case Study~1, with one charge per gradient ($c_g=3$; Section~\ref{S-sec:S-xp-grad}), and 220 of their 240 strict feasibility labels cannot be decided from the stored coordinates (Section~\ref{sec:archive-reliability}), so their IEA37 lead rests on the labels reported by the runs (worst case: Table~\ref{S-tab:S-sa-feasworst});
\item turbulence and loads (critical at $4D$), wake steering, blockage, terrain, cables, electrical losses and costs are outside our scope; all AEP values are gross.
\end{itemize}
% CHECK-DIR [F03]-[F07]: coarse bins exploited; [F21]: 0 of 961 feasible runs above the installed block at 1 deg; [F32]/[F34]: PyWake 661.39 vs ours 666.75 GWh/yr, wake loss -0.72 pp, cause C_T at free-stream speed; [F44]: projected IEA37 gaps -1.93 / -6.23 % (30,030)
% CHECK-FINAL [C45] (budget study; NBudCloseGapOneTwentyK generated); CHECK-EXTRA [X52]: Horns Rev feasibility (one site, random starts); CHECK (manual, 2026-10-04): GA lead at 30,030 under 1-deg bins and PyWake NOJ (rev3_sites.json); Lillgrund PSO-VNS lead over the runners-up SSA-VNS, VNS, MS-SLSQP n.s. at 1 deg (p_Holm 0.16, 0.58, 0.58); still significant vs PSO, SSA, LX-SSA, DE, RSD-VNS

\emph{Statistical conclusion validity.}
\begin{itemize}
\item The margin, the $N\ge10$ split, the Data Set~II subgroup and the density pattern were stated after the primary analysis; the margin is a convention, and the minimal margins let readers apply their own.
\item The 68 cases form only six clusters with nested $N$: an exact sign test over clusters cannot fall below $p=0.031$, cluster-robust intervals have 5 degrees of freedom, and the p-values describe this benchmark, not a population of farms.
\item Seeds are synchronized across cases; resampling them jointly leaves the PSO-VNS vs.\ PSO verdicts unchanged, but not the seed-level equivalence of SSA-VNS and RSD-VNS (Table~\ref{tab:equiv-main}).
\item The all-run sign tests treat case--seed pairs as independent; with seeds as the unit of replication, no all-run verdict changes (Section~\ref{S-sec:S-sa-dep}).
\item The budget and feasible-start rankings (six cases) are not tested, and the feasibility result rests on two sites with 30 seeds each.
\item The Lillgrund, spacing, GA, exact-gradient, constraint-handling and direct 1$^\circ$ studies were conducted after the primary analysis, the last two on six and eight cases; equivalence statements other than the two primary questions (Section~\ref{sec:setup-stats}) are exploratory, and their error rates are not controlled jointly.
\end{itemize}
% CHECK-FINAL [C49] + CHECK-EXTRA [X41]: post hoc margin, minimal margins; [X47]: density pattern (exploratory); [X10]: six clusters; [X11]: exact p = \NXClMinP; [X42]: CR2 5 d.f.; CHECK-FINAL [C29]: rank differences not tested; CHECK-EXTRA [X52]; CHECK (manual, 2026-10-04): joint seed resampling (tab:S-sa-syncseed): PSO-VNS vs PSO [-0.026, -0.010] unchanged; SSA-VNS vs RSD-VNS [-0.002, 0.050], seed-level equivalence not robust, non-superiority robust

\subsection{Conclusions and Future Directions}
\label{sec:conclusion}
We proposed PSO-VNS, a two-phase method for the WFLOP that runs the standard PSO with constriction coefficients for half of the evaluation budget and then a variable neighborhood search from the best layout found. We compared it with eight methods (PSO, SSA-VNS, SSA, LX-SSA, DE, GA, VNS and MS-SLSQP) on four types of wind farm, with equal evaluation budgets, seed-paired runs, random starts and the implemented penalty and box repair. The findings are:
\begin{enumerate}
\item \emph{Type~I, Kusiak--Song benchmark (68 cases).} PSO-VNS has the best average rank; it is not significantly different from PSO and GA and significantly better than the other six methods. Its average advantage over PSO is small under the 15$^\circ$ rose (within the post hoc margin on this fixed benchmark and under exchangeable-case sensitivity, not under scenario-cluster sensitivity), with case-level differences in both directions; on the separate all-run endpoint, PSO-VNS is better more often than not, also with seeds as the unit of replication.
\item \emph{Type~II, Horns Rev~1 block.} PSO-VNS ranks first at 6,030 evaluations, feasible in all 30 runs (PSO: 19); at 30,030 evaluations GA ranks first and PSO-VNS second.
\item \emph{Type~III, Lillgrund block.} PSO-VNS ranks first by mean AEP at both budgets, but at 6,030 evaluations MS-SLSQP is more reliable (feasible in all 30 runs, PSO-VNS in 17, PSO in none); the AEP lead of PSO-VNS over the runners-up is not significant with 1$^\circ$ bins.
\item \emph{Type~IV, IEA37 Case Study~1.} PSO-VNS ranks first with 16 turbines and GA with 36 turbines, at both budgets. On this piecewise-smooth model, SLSQP with exact gradients, alone or after a PSO phase, does better than every metaheuristic compared, with feasibility labels as reported by the runs (220 of its 240 strict labels cannot be decided from the stored coordinates).
\item \emph{Ablation.} The PSO phase is the component that matters: it gives the second phase, essentially one compass-search descent, a clearly better start than random sampling in the farm, whereas at the benchmark spacing $4D$ salp-swarm first phases give no gain larger than the 0.05~pp margin over random sampling in the farm (ruled out across seeds and cases; at $6D$ the SSA phase does gain).
\item \emph{Robustness.} On the benchmark, PSO-VNS keeps first place when its layouts are re-evaluated with 1$^\circ$ direction bins and when the search itself uses 1$^\circ$ bins (eight cases). With direct 1$^\circ$ optimization, its advantage over PSO exceeds the margin and is about twice that in a seed-matched control arm optimized with 15$^\circ$ bins and evaluated with 1$^\circ$ bins (difference-in-differences $-0.115$~pp, 95\% CI [$-0.167$, $-0.060$]). The comparison depends partly on constraint handling (six cases, five methods): with Deb's feasibility rules on normalized violations and box clipping, the order of the leading methods changed; with radial projection, the run-level order and the leading method persisted, but PSO-VNS's all-run advantage over PSO was not significant.
\end{enumerate}
% CHECK-FINAL [C53]/[C54]: no SSA gain over disc sampling = "in the farm"; [C55]: LX-SSA-VNS worse; [C56]/[C14]/[C44]: PSO phase beats both controls; CHECK-EXTRA [X53]: gain > margin ruled out across seeds and cases (one-sided); CHECK-DIAG [D08]: essentially one descent
% CHECK-DIR [F06]/[F09] + CHECK-FINAL [C21]/[C22]: PSO-VNS first at 15/5/1 deg (Jensen), cubic curves, Gaussian 1 deg; PSO first only under Gaussian 15 deg; CHECK-EXTRA [X39]: seed/case equivalent, cluster not; CHECK-EXTRA [X43]: 19/68 beyond margin, both directions; CHECK (manual, 2026-10-04): all-run score 0.551 (tab:S-sa-allrun); direct 1-deg: PSO-VNS rank 1.12, PSO-VNS - PSO -0.194 pp, case 90% [-0.272, -0.114] beyond -0.05 (rev3_fine.json)
% CHECK-EXTRA [X52]: Horns Rev 30/30 vs 19/30, McNemar p 9.8e-4; CHECK-FINAL [C08]/[C20]; CHECK (manual, 2026-10-04): Lillgrund 17/30 vs 0/30 and 30/30 vs 19/30; GA first at 30,030 under 5-deg, 1-deg and PyWake NOJ (tab:S-sa-hr-pool, tab:S-sa-hr-eval); Lillgrund 1 deg p_Holm 0.16 / 0.58 vs SSA-VNS / MS-SLSQP and VNS (tab:S-sa-lg-eval); constraint study, 6 cases x 5 methods, projection: conditional dL -0.095 pp [-0.162, -0.037], all-run 0.561 ns (tab:S-sa-constraint)

\emph{Future directions.} Future work follows from the threats above:
\begin{itemize}
\item \emph{Expensive wake models}: surrogate-assisted evaluation (graph neural network, data-driven or Kriging surrogates~\cite{Yang2023,Wang2024,Li2025}) would make equal-budget comparisons at higher fidelity affordable, if the surrogate error is reported next to the method differences.
\item \emph{Larger farms and more sites}, such as the IEA37 64-turbine scenario and the full Horns Rev~1 and Lillgrund farms.
\item \emph{Gradient-based baselines} on smooth models of larger farms~\cite{Rodrigues2024}, and swarm--gradient hybrids, under budgets that charge every gradient.
\item \emph{Stronger baselines and constraint handling}: CMA-ES, L-SHADE, tuned DE, and constraint handling designed for each method family.
\item \emph{Fine roses during optimization} on all cases, with loads and economic objectives.
\end{itemize}
